% Horizon first law: radial weight yes; CHM parabola not unique.
\documentclass[11pt]{article}
\usepackage[margin=1.05in]{geometry}
\usepackage{amsmath,amssymb,amsthm}
\usepackage[colorlinks=true,linkcolor=blue,citecolor=blue,urlcolor=blue]{hyperref}

\newtheorem*{admission}{Admission}

\title{\textbf{On the Cut, a Radial Weight Beats a Flat Kernel;\\
the CHM Parabola Is Not Selected}}
\author{Robert W.\ McGwier\\
\small CTO, Cohere Technology Group\\
\small GitHub: \texttt{n4hy}}
\date{15 August 2026 --- Version 1}

\begin{document}
\maketitle

\begin{abstract}
\noindent
Paper~27 found that on the horizon, $K_{\mathrm{CHM}}$ beats a
flat hop kernel. Paper~25 found that for the modular kernel
itself, a linear $R-r$ weight can beat the CHM parabola on a
small ball. This note asks the first-law form of that question,
on a new grid, discarding any hop that flips $n_{\mathrm{occ}}$.

$N=12$, $R=4$, $294$ surface hops, zero occupancy flips.
$R_{\mathrm{CHM}}=0.808<R_{\mathrm{flat}}=0.994$ (C3 PASS).
$R_{\mathrm{CHM}}=0.8082$ versus $R_{\mathrm{lin}}=0.8099$
(C4 is a tie). Tracking floors fail
($\lvert\rho\rvert=0.589<0.60$).

An auditor on $N=14$, $R=3$ kept only $8$ hops after $142$
occupancy flips and is not a C4 confirmation.

The cut prefers a radially decreasing weight to a shapeless
local kernel. It does not prefer $R^2-r^2$ to $R-r$.
Not $\eta=1/4G$. Not Einstein. Not the Bloch guess.
\end{abstract}

\section{Numbers}

\begin{center}
\begin{tabular}{lcccc}
\hline
 & $n$ & $R_{\mathrm{CHM}}$ & $R_{\mathrm{lin}}$ & $R_{\mathrm{flat}}$ \\
\hline
$N=12,R=4$ & $294$ & $0.808$ & $0.810$ & $0.994$ \\
audit $N=14,R=3$ & $8$ & $0.959$ & $0.959$ & $0.982$ \\
\hline
\end{tabular}
\end{center}

C3 is the Paper~27 fact, now at a second radius with a stable
Fermi sea. C4 is the new question and it does not select the
quadratic.

\begin{admission}
I asked whether the parabola was doing Clausius work beyond any
decreasing radial weight. On this grid it is not. I did not
call a $0.002$ residual gap a theorem, and I did not write
Einstein.
\end{admission}

\end{document}
